% content/c05-teilchen/en.tex — Chapter 5: Which Hilbert space? The particle on a line
% (English version; German leads)
%
% Status 2026-11-11: a beginning. One Kern section carries the two figures of lectures 7 and 8
% and their demos (docs/plan/lec78/DEMOS.md). Prose and formulas from script v0.3:
% latex/lectures/en/lec07.tex and lec08.tex. Notation as in lecture 8.

\chapter{Which Hilbert space? The particle on a line}\label{c05}

% Voice: lec07.tex.
The postulates of Chapter~\ref{c04} do not say which Hilbert space, which Hamiltonian and which
observables belong to a given system. For the most important system, the particle on a line,
the Hilbert space is $\Hil = \Ltwo{\mathbb{R}}$, the space of square-integrable functions
$\Psi\colon \mathbb{R} \to \mathbb{C}$ with the scalar product
$\braket{\Phi}{\Psi} = \int \Phi^*(x)\Psi(x)\,dx$. One writes the value of the function at the
point $x$ as $\Psi(x) = \braket{x}{\Psi}$ and calculates with the kets $\ket{x}$ as if they were
an orthonormal basis, the \emph{position basis}; the Kronecker delta is replaced by the Dirac
delta function, $\braket{x}{y} = \delta(x-y)$, and sums over the basis become integrals. The
function $\Psi$ is the list of coefficients of $\ket{\Psi}$ in the position basis; it is called
the \emph{wave function}.

\section{Position density and the free particle}\label{c05.s1}

A position detector monitors an interval $[a,b]$ and answers “yes” if it finds the particle
there. Its operator is
\begin{equation}\label{c05.eq.ortsdetektor}
  \op{\EffOp}_{[a,b]} = \int_a^b \ketbra{x}{x}\,dx .
\end{equation}
It responds with probability
\begin{equation}\label{c05.eq.intervall}
  \Prob\bigl[x\in[a,b]\bigr] = \melem{\Psi}{\op{\EffOp}_{[a,b]}}{\Psi}
  = \int_a^b \abs{\Psi(x)}^2\,dx
\end{equation}
in the pure state $\ket{\Psi}$. Since $\abs{\Psi(x)}^2$ is not negative, $p(x) = \abs{\Psi(x)}^2$ is a probability density, the
density of the position measurement. The wave function itself assigns a \emph{complex} number
to every position; you should resist the temptation to picture $\Psi$ itself as a density.

\begin{tpcfigure}{c05.fig.dichte}
  \demo{c05.demo.detektor}
  \caption{The same wave function twice: the Gaussian wave packet
    $\Psi(x) = (2\pi)^{-1/4} e^{-x^2/4} e^{3ix}$. (a) Real and imaginary part change sign all
    the time; dotted, the envelope $\pm\abs{\Psi}$. (b) Of the factor $e^{3ix}$ nothing is left
    in $p = \abs{\Psi}^2$. The area over $[a,b]$ is the response probability of the detector
    $\op{\EffOp}_{[a,b]}$, Equation~\eqref{c05.eq.intervall}. The narrow strip of width
    $\varepsilon$ at $c$ is the same construction, only narrower: as $\varepsilon$ shrinks, the
    area shrinks with it.}
\end{tpcfigure}

Why an interval and not a point? The probability of finding the particle \emph{exactly} at the
point $x$ is always zero. The object $\ketbra{x}{x}$, taken on its own, is not an operator:
$(\ketbra{x}{x})^2 = \delta(0)\ketbra{x}{x}$, and $\delta(0)$ is not a number. For an interval
the delta disposes of one of the two integrals, and
$\op{\EffOp}_{[a,b]}^2 = \op{\EffOp}_{[a,b]}$ is a projector.

How does a free particle of mass $m$ move? Its Schrödinger equation
$i\hbar\,\partial_t\Psi = -\frac{\hbar^2}{2m}\,\partial_x^2\Psi$ is solved by every plane wave
$e^{i(kx - \hbar k^2 t/2m)}$, and since it is linear, also by every superposition:
\begin{equation}\label{c05.eq.freieloesung}
  \Psi(x,t) = \frac{1}{\sqrt{2\pi}}\int_{-\infty}^{\infty} \widetilde{\Psi}(k)\,
  e^{i\left(kx-\frac{\hbar k^2}{2m}t\right)}\,dk .
\end{equation}
The weights $\widetilde{\Psi}(k)$ are fixed by the initial condition: at $t = 0$ this is the
inverse Fourier transformation, and $\widetilde{\Psi}$ is the Fourier transform of the initial
wave function,
\begin{equation}\label{c05.eq.phik}
  \widetilde{\Psi}(k) = \frac{1}{\sqrt{2\pi}}\int_{-\infty}^{\infty} \Psi(x,0)\,e^{-ikx}\,dx .
\end{equation}

\begin{merksatz}[The free particle in three steps]
  Compute the Fourier transform $\widetilde{\Psi}(k)$ of the initial wave function
  $\Psi(x,0)$. Multiply each component $\widetilde{\Psi}(k)$ by its phase
  $e^{-i\hbar k^2t/2m}$. Transform back.
\end{merksatz}

For the Gaussian wave packet
$\Psi(x,0) = (\pwidth\sqrt{2\pi})^{-1/2}\, e^{-(x-a)^2/4\pwidth^2}$ the forward and backward
transformations are Gaussian integrals. The squared modulus remains a bell curve about $a$,
whose width grows:
\begin{align}
  \abs{\Psi(x,t)}^2 &= \frac{1}{\pwidth_t\sqrt{2\pi}}\;
  e^{-\frac{(x-a)^2}{2\pwidth_t^2}}, \label{c05.eq.zerfliessen}\\
  \pwidth_t &= \pwidth\sqrt{1+\frac{t^2}{\tspread^2}},
  \qquad
  \tspread = \frac{2m\pwidth^2}{\hbar}. \notag
\end{align}
The wave packet \emph{spreads}: after the time $\tspread$ it is broader by a factor $\sqrt{2}$,
and the narrower it was initially, the shorter $\tspread$ is.

The same information is contained in the \emph{momentum wave function} $\Phi(p,t)$, the
coefficients of $\ket{\Psi}$ in the basis of plane waves with momentum $p = \hbar k$; at time
zero $\Phi(p,0) = \widetilde{\Psi}(p/\hbar)/\sqrt{\hbar}$. For the free particle only a phase
changes in momentum space,
\begin{equation}\label{c05.eq.impulszeit}
  \Phi(p,t) = e^{-\frac{it}{\hbar}\frac{p^2}{2m}}\,\Phi(p,0),
  \qquad
  \abs{\Phi(p,t)}^2 = \abs{\Phi(p,0)}^2 .
\end{equation}
For the Gaussian wave packet of width $\pwidth$, $\abs{\Phi(p,0)}^2$ is a Gaussian
distribution about $p = 0$ with standard deviation $\pwidth_p = \hbar/2\pwidth$; at time zero
the product of the two widths is $\pwidth\,\pwidth_p = \hbar/2$.

\begin{tpcfigure}{c05.fig.fourier}
  \demo{c05.demo.zerfliessen}
  \caption{The Gaussian wave packet at time $t = 0$, three times, with widths $\pwidth/2$,
    $\pwidth$ and $2\pwidth$, on the left in position space, on the right in momentum space;
    the same line means the same packet. The bars above the curves mark $\pm$ one standard
    deviation: halving the width in position space doubles it in momentum space. The order of
    the widths is reversed, the product stays. In time the curves on the left spread according
    to Equation~\eqref{c05.eq.zerfliessen}, those on the right stay put according to
    Equation~\eqref{c05.eq.impulszeit}.}
\end{tpcfigure}
